\documentclass[11pt,a4paper]{article}
\usepackage[utf8]{inputenc}
\usepackage[T1]{fontenc}
\usepackage[french]{babel}
\usepackage{geometry}
\usepackage{xcolor}
\usepackage{hyperref}
\usepackage{titlesec}

\geometry{margin=1in}
\definecolor{titleblue}{RGB}{0,51,102}

\hypersetup{
    colorlinks=true,
    linkcolor=titleblue,
    urlcolor=titleblue,
    pdftitle={Lettre de Motivation - Ismail AISSAOUI},
}

\begin{document}

\noindent \textbf{Ismail AISSAOUI} \\
Ingénieur d'État en Intelligence Artificielle (ESI-SBA) \\
Email: ismail.aissaoui.pro@gmail.com \\
Téléphone: +213 660 70 77 96 \\
LinkedIn: https://ismail-aissaoui.vercel.app

\bigskip

\begin{flushright}
    Au Comité de Sélection du Doctorat \\
    \textbf{Programme EDT - Université de Lille / INSA Rennes} \\
    À l'attention de : Prof. Romain Rouvoy et Dr. Quentin Perez \\
    \medskip
    1er mai 2026
\end{flushright}

\bigskip

\noindent \textbf{Objet : Candidature au doctorat : Améliorer la durabilité dans le continuum physique-numérique des jumeaux (Rang 8 - PC3)}

\bigskip

Monsieur le Professeur, Monsieur le Docteur,

Ingénieur d'État en Intelligence Artificielle diplômé de l'ESI-SBA, je vous adresse avec enthousiasme ma candidature pour le poste de doctorat intitulé « Améliorer la durabilité dans le continuum physique-numérique des jumeaux » au sein du programme national Engineering Digital Twin (EDT). Cette candidature est motivée par la pertinence de mon expertise en déploiement d'IA sur matériel contraint pour vos recherches sur le middleware durable.

Actuellement en fin de cycle ingénieur chez Sonatrach, je travaille sur le développement d'un **Jumeau Numérique informé par la physique** pour la surveillance de turbines à gaz. Un volet crucial de mon travail consiste à optimiser le déploiement de modèles hybrides haute performance (**Mamba-SSM**, **xLSTM**) sur des dispositifs industriels en périphérie (Edge). Cette expérience m'a confronté au compromis direct entre précision du modèle et consommation énergétique. J'ai consacré une partie importante de mes recherches à l'analyse de l'empreinte énergétique de mes modèles pour garantir un fonctionnement durable sur des infrastructures isolées.

Je suis particulièrement attiré par votre projet axé sur le **déploiement adaptatif** et l'analyse du cycle de vie des jumeaux numériques. Ma maîtrise de l'IA distribuée et de l'optimisation logicielle s'aligne parfaitement avec votre objectif de créer un framework de middleware capable de sélectionner dynamiquement la fidélité du modèle et la plateforme d'exécution (Cloud vs Edge) en fonction de contraintes énergétiques.

Rejoindre l'Université de Lille et l'INSA Rennes au sein du consortium EDT représente pour moi une opportunité unique de contribuer à l'émergence d'un numérique plus sobre et responsable au sein de la **plateforme Artemis**.

Je vous remercie de l'attention que vous porterez à ma candidature et reste à votre entière disposition pour un entretien.

Dans l'attente de votre réponse, je vous prie d'agréer, Monsieur le Professeur, Monsieur le Docteur, l'expression de mes salutations distinguées.

\bigskip

\noindent Ismail AISSAOUI

\end{document}
